% =================================================================
% GLOSSARY
% =================================================================
% This file contains a consolidated list of key terms from the book.
% =================================================================

\chapter{Glossary}

\begin{description}

    \item[Antonym] A word with the opposite meaning of another.

    \item[Apostrophe] Mark showing a missing letter or ownership.

    \item[Audience] The person or people meant to read a piece of writing.

    \item[Character] Who a story is about.

    \item[Comma] Mark that separates items or adds a pause.

    \item[Comprehension question] A question about what a passage means.

    \item[Context] The situation or sentence a word is used in.

    \item[Context clue] A hint in the surrounding text that helps explain a word's meaning.

    \item[Contraction] Shortened form of two words (e.g., don't).

    \item[Diary entry] A first-person account of a day's events, usually in past tense.

    \item[Editing] Checking and correcting a piece of writing.

    \item[Informal letter] A personal letter written to someone you know well.

    \item[Issuing authority] The person or role responsible for a notice (e.g., Principal).

    \item[Language-focus question] A question about a specific grammar or vocabulary feature.

    \item[Language workshop] Using one passage for several different language tasks.

    \item[Message] A short note giving essential information quickly.

    \item[Narrative] A piece of writing that tells a story.

    \item[Notice] A short, official piece of writing informing a group.

    \item[One-word expression] A single word that replaces a longer phrase.

    \item[Paragraph] A group of sentences developing one main idea.

    \item[Possession] Showing that something belongs to someone.

    \item[Prefix] A word part added to the front to change meaning.

    \item[Punctuation] Marks that make written meaning clear.

    \item[Quotation marks] Marks that enclose a speaker's exact words.

    \item[Root word] The main word that carries the core meaning.

    \item[Sensory detail] A specific detail about sight, sound, smell, touch, or taste.

    \item[Sequencing word] A word (first, then, finally) showing the order of events.

    \item[Setting] Where and when a story happens.

    \item[Spatial order] Arranging details by location in space.

    \item[Suffix] A word part added to the end to change form or meaning.

    \item[Supporting detail] A sentence that develops or explains the topic sentence.

    \item[Synonym] A word with almost the same meaning as another.

    \item[Time order] Arranging details in the order they happened.

    \item[Topic sentence] The sentence that states the paragraph's main idea.

    \item[Transformation] Rewriting a sentence in a different grammatical form without changing its meaning.

    \item[Unity] Every sentence in a paragraph relating to the same main idea.

    \item[Word bank] A collected list of useful words.

    \item[Word family] A group of words sharing the same root.

\end{description}